\documentclass{article}
\usepackage{oeneye-font}

\begin{document}

\begin{center}
    {\Huge \textoeneye{oeneyeOS Kernel Specifications}} \\[1em]
    {\Large \textoeneye{Subsystem Architecture \& Operator Framework}} \\[2em]
    \textbf{Author:} Kai Olaf Ketelhut \\
    \textbf{Target Architecture:} x86 / Emulated Legacy Runtime \\
    \today
\end{center}

\hrule
\vspace{1.5em}

\section*{\textoeneye{1. Overview}}
The \textoeneye{oeneyeOS} runtime architecture couples low-level binary image structures with deterministic operator pipelines. By utilizing extracted raw bitmap font containers (`OMQ.FNT`) transformed into scalable vector fonts, all diagnostic readouts and low-level subsystem logs maintain pixel-exact fidelity across documentation targets.

\section*{\textoeneye{2. Core Subsystems}}
\begin{itemize}
    \item \textoeneye{Memory Manager}: Direct segment mapping and localized allocation tracking.
    \item \textoeneye{Operator Kernel}: Symbolic execution loop running autopoietic routines.
    \item \textoeneye{Rendering Pipeline}: Vector-mapped bitmap font integration via XeLaTeX.
\end{itemize}

\end{document}
